\chapter{源文件与默认参数}

\section{源文件}\label{app:files}

RTL 共 \RtlFileCount{} 个 Verilog 文件（合计 \countsep{\RtlLineCount} 行，\DesignUnitCount{} 个设计单元），均位于仓库 \code{rtl/} 目录。

\begin{longtable}{>{\raggedright\arraybackslash}p{4.6cm}>{\raggedright\arraybackslash}p{9.4cm}}
\caption{RTL 源文件清单（\RtlFileCount{} 个文件）}\label{tab:files}\\
\longtablehead{2}{文件 & 作用}
\code{AFDX_top.v} & 赛题模板顶层封装：2B VLID 流式包装器、配置窥探与 \code{afdx_es} 核心 \\
\code{afdx_es.v} & 端系统核心顶层：双网收发核心、服务通道、三项服务、配置与时间基准 \\
\code{afdx_app_tx.v} & 顶层发送流包装器：2B VLID 前缀解析、3 级 128 路比较映射与 32~KiB 缓存 \\
\code{afdx_app_rx.v} & 顶层接收流包装器：消息缓存、VLID 反查、2B VLID 前缀插入与溢出保护 \\
\code{afdx_app_map.v} & 接收包装层 $4096\times 16$~bit $\text{idx}\rightarrow\text{VLID}$ 专用同步映射 RAM \\
\code{afdx_app_payload.v} & 包装层 $8\times 4$~KiB 分体同步载荷存储（含 \code{_bank}/\code{_slice} 及写隔离） \\
\code{afdx_app_read.v} & 包装层 8 信用流式同步预取队列（屏蔽 4 拍块 RAM 读延迟） \\
\code{afdx_timebase.v} & 周期计数 \code{now}、粗粒度计时 \code{now_c} 与字节使能 \code{byte_en} \\
\code{afdx_cfg.v} & 18 位字地址译码、寄存器、27 个计数器两级流水维护、配置表写握手 \\
\code{afdx_svc.v} & 应用与服务的整消息轮询发送（授权寄存）、按 UDP 目的端口分发接收 \\
\code{afdx_tx_core.v} & 发送核心集成模块，例化发送各子模块、输出流缓冲及 \code{afdx_tx_payload_store} \\
\code{afdx_tx_buf.v} & $1\text{--}\MaxMessageBytes$~B 分页写入、两槽入口弹性级与 $16\times 8$ 回滚 RAM \\
\code{afdx_tx_free.v} & 256 页空闲页 FIFO，支持双页面回收端口并发回收 \\
\code{afdx_tx_engine.v} & BAG 与分片调度：扫描、入队、单帧/分片授权、分片进度表维护与起始确认 \\
\code{afdx_tx_builder.v} & 42B/34B 组帧、12 拍 IPv4 校验和准备、载荷预读与 $-8$ 拍提前补偿 \\
\code{afdx_mac_tx.v} & 双网发送 MAC：前导码、FCS、帧间隔、发送下溢（TX Underrun）检测与异常终止 \\
\code{afdx_crc32_d8.v} & 逐字节 8 位并行 CRC-32 组合逻辑 \\
\code{afdx_rx_core.v} & 接收核心集成模块，例化双 MAC、解析器、32-Bank 双写分片映射表、重组与交付 \\
\code{afdx_mac_rx.v} & 接收 MAC：SFD 检测、CRC 校验、剥离 FCS \\
\code{afdx_rx_parse.v} & 单帧/首片与后续分片帧头检查、载荷预写入环形缓冲、帧尾分类 \\
\code{afdx_rx_engine.v} & 30 拍日程：32-Bank 查表、配置检查、IC、RM、重组槽位查询预约与过期清理 \\
\code{afdx_rx_deliver.v} & 读取 $32\times 73$~bit \code{rx_dq}，从双网环形缓冲逐字节流式送出 \\
\code{afdx_rx_reassembly.v} & 双 Bank $\times 8$ 槽位（\ReasmContexts{} 上下文）顺序 IPv4 分片重组顶层 \\
\code{afdx_rx_reassembly_bank.v} & 单 Bank 8 槽位重组状态机、连续偏移检查、整包 UDP 校验与 2~s 超时回收 \\
\code{afdx_rx_reassembly_payload.v} & $16\times 8$~KiB 重组载荷存储池（64 条 $8192\times 2$ 通道与 128 个写隔离 LUT） \\
\code{afdx_rx_reassembly_ram.v} & 分片重组元数据与接收窄分片表复用的专用同步双端口 RAM 封装 \\
\code{afdx_ram_sdp.v} & 通用简单双端口（SDP）同步 RAM 原语封装（例化 \RamSdpCount{} 处） \\
\code{afdx_rtc.v} & 实时时钟与 SNTPv4 纳秒级时钟同步（含 40~ns 流水补偿） \\
\code{afdx_ntp_fraction.v} & 纳秒与 NTP 32 位二进制小数的 32 轮逐位换算 \\
\code{afdx_snmp.v} & SNMP 顶层包装：请求捕获、协议 RAM 接口与 \code{CE_DIV=2} 使能树 \\
\code{afdx_snmp_engine.v} & SNMP 协议核心：ASN.1/BER 编解码、MIB 读写与响应生成（16~ns 域） \\
\code{afdx_615a.v} & 615A 顶层包装：TFTP 捕获、6 块协议 RAM 接口与 \code{CE_DIV=2} 使能树 \\
\code{afdx_615a_engine.v} & 615A 协议核心：数据加载/卸载状态机与文件 CRC-16 硬件校验（16~ns 域） \\
\code{afdx_mgmt_tx.v} & 管理面流式预取发送适配器：8 字节预取队列，供 SNMP 与 615A 复用解耦 \\
\end{longtable}

\section{默认参数}\label{app:param}

\begin{longtable}{>{\raggedright\arraybackslash}p{4.0cm}>{\raggedleft\arraybackslash}p{2.5cm}>{\raggedright\arraybackslash}p{\dimexpr\linewidth-6.5cm-6\tabcolsep\relax}}
\caption{顶层与核心参数及默认容量}\label{tab:param}\\
\longtablehead{3}{参数 & 默认值 & 含义}
\code{TX_PAGE_AW} / \code{TX_OFF_AW} & 8 / 9 & 发送缓存 256 页 $\times 512$~B，共 128~KiB \\
\code{LOOKAHEAD} & 512 & 提前选择窗口，512 拍（4.096~µs） \\
\code{RX_RING_AW} / \code{RX_DQ_AW} & 12 / 5 & 每网接收环形缓冲 4~KiB，交付队列 32 项（73 位宽） \\
\code{RX_CONFIG_BANKS} & 16 & 接收配置总线与描述符局部复制扇出分支数 \\
\code{CE_DIV} & 2 & SNMP / 615A 协议引擎多周期时钟使能分频比（16~ns 域） \\
\code{REASSEMBLY_CONTEXTS} & \ReasmContexts & 接收 IPv4 顺序分片重组并发槽位数（2 Bank $\times 8$ 槽 $\times 8$~KiB） \\
\code{REASSEMBLY_TIMEOUT} & $2.5\times 10^8$ 拍 & 分片重组超时老化阈值（2.0~s） \\
\midrule
单条应用消息长度 & $1\text{--}\MaxMessageBytes$~B & $>1471$~B（或 $>\text{max\_payload}$）自动按 IPv4 分片传输 \\
顶层应用包装缓存 & $2\times 32$~KiB & \code{afdx_app_tx} 与 \code{afdx_app_rx} 各 32~KiB + 32 项描述符 \\
发送 / 接收 VL 容量 & 128 / 4096 & 接收映射表 65536 项；声明满载规划为 63 条并发收发 VL \\
帧间隔（IFG） & 12 拍 & 96~ns \\
上电表清零初始化 & 65,539 拍 & 核心 \code{init_done}（524.312~µs）；顶层 \code{tx_ready} 约 65,550 拍 \\
SNMP 请求缓存 & 256~B & 单 VarBind，支持 36 个 MIB 标量/表格节点 \\
615A 加载区 & 8~KiB & 头文件区 2~KiB，LUR 区 512~B，超时 2~s / 重传 1 次 \\
RTC 同步周期 & 125000 拍 & 1~ms，可经寄存器修改 \\
\end{longtable}

\chapter{服务寄存器}\label{app:regs}

基本寄存器、计数器和配置表地址见表~\ref{tab:addrmap}，服务寄存器见表~\ref{tab:regs}。RW 为可读写，RO 为只读。

\begingroup
\renewcommand{\tablestyle}{\small\linespread{1.04}\selectfont
  \renewcommand{\arraystretch}{1.10}}
\begin{longtable}{>{\raggedright\arraybackslash}p{2.4cm}>{\raggedright\arraybackslash}p{3.6cm}>{\centering\arraybackslash}p{0.9cm}>{\raggedright\arraybackslash}p{7.0cm}}
\caption{服务寄存器}\label{tab:regs}\\
\longtablehead{4}{地址 & 名称 & 读写 & 含义}
\multicolumn{4}{l}{\textit{RTC}} \\*
\code{0x80} & \code{CTRL} & RW & 位 0 启用，位 1 Server 模式；写位 2 装入暂存时间 \\
\code{0x81} & \code{PERIOD} & RW & Client 端请求间隔（拍），复位 125000 \\
\code{0x82} & \code{TX_VL} & RW & 发送 VL，低 7 位 \\
\code{0x83} & \code{UDP_PORT} & RW & UDP 端口，复位 123 \\
\code{0x84} & \code{SEC} & RW & 写秒暂存值；读当前秒并锁存纳秒 \\
\code{0x85} & \code{NS} & RW & 写纳秒暂存值；读锁存的纳秒值 \\
\code{0x86} & \code{OFFSET_SEC} & RO & 最近偏差的有符号秒部分 \\
\code{0x87} & \code{OFFSET_NS} & RO & 最近偏差的纳秒余数 \\
\code{0x88} & \code{SYNC_COUNT} & RO & Client 时钟同步次数 \\
\midrule
\multicolumn{4}{l}{\textit{SNMP}} \\*
\code{0x90} & \code{ENABLE} & RW & 位 0 启用代理 \\
\code{0x91} & \code{TX_VL} & RW & 响应与 Trap 的发送 VL \\
\code{0x92} & \code{COMMUNITY_LENGTH} & RW & 1--16，复位 6 \\
\code{0x93}--\code{0x96} & \code{COMMUNITY_0}--\code{3} & RW & community 内容，复位为 \code{public} \\
\code{0x97} & \code{SAMPLE_CYCLE} & RO & 上次 MIB 读取时的 \code{now} \\
\code{0x98} & \code{UPTIME} & RO & sysUpTime，单位 10~ms \\
\code{0x99} & \code{REQUESTS} & RO & 端口匹配的请求数 \\
\code{0x9A} & \code{BAD_COMMUNITY} & RO & community 不符次数 \\
\code{0x9B} & \code{TRAP_ENABLE} & RW & 位 0 启用 Trap \\
\code{0x9C} & \code{TRAPS_SENT} & RO & 已发送 Trap 数 \\
\code{0xA0}--\code{0xA3} & \code{DESCRIPTION_0}--\code{3} & RO & sysDescr 字符串 \\
\midrule
\multicolumn{4}{l}{\textit{615A}} \\*
\code{0xB0} & \code{ENABLE} & RW & 位 0 启用 \\
\code{0xB1} & \code{TX_VL} & RW & 发送 VL \\
\code{0xB2} & \code{STATUS} & RO & 位 1:0 空闲/进行中/成功/失败；位 15:8 原因码 \\
\code{0xB3} & \code{LOADED_BYTES} & RO & 本次已写字节数 \\
\code{0xB4} & \code{MEMORY_ADDRESS} & RW & 加载区读回地址，低 13 位 \\
\code{0xB5} & \code{MEMORY_DATA} & RO & 加载区读回字节 \\
\code{0xB6} & \code{LUS_COUNTER} & RO & 本次状态文件数 \\
\code{0xB7} & \code{RETRY_TOTAL} & RO & 累计超时重传次数 \\
\code{0xB8} & \code{REQUEST_PORT} & RW & TFTP 请求端口，复位 59 \\
\code{0xB9} & \code{TIMEOUT_CYCLES} & RW & 超时拍数，复位 250000000 \\
\code{0xBA} & \code{RETRY_LIMIT} & RW & 重传上限，复位 1 \\
\code{0xBB} & \code{PROTOCOL_VERSION} & RW & \code{A4} 或 \code{A3} \\
\end{longtable}
\endgroup

\chapter{内部接口信号}\label{app:intif}

本附录按信号说明第~\ref{sec:intif}~节汇总的内部接口。

\section{消息接口}

服务通道 \code{afdx_svc} 把选中的消息以字节流交给发送核心，见表~\ref{tab:ifm}。没有服务报文时，这组信号与应用口组合直通。

\begin{longtable}{>{\raggedright\arraybackslash}p{4.4cm}>{\raggedleft\arraybackslash}p{1.8cm}>{\raggedright\arraybackslash}p{\dimexpr\linewidth-6.2cm-6\tabcolsep\relax}}
\caption{消息接口（\texttt{afdx\_svc} $\rightarrow$ \texttt{afdx\_tx\_buf}）}\label{tab:ifm}\\
\longtablehead{3}{信号 & 位宽 & 含义}
\code{m_valid}、\code{m_ready} & 1、1 & 字节握手；\code{m_ready} 只在 \code{tx_buf} 处于 IDLE、RECV、DRAIN 时为 1 \\
\code{m_data}、\code{m_last} & 8、1 & 消息字节、末字节标记 \\
\code{m_vl} & 7 & 发送 VL 索引，首字节锁存 \\
\code{m_port_en}、\code{m_sport}、\code{m_dport} & 1、16、16 & \code{m_port_en} 为 1 时使用消息自带的 UDP 端口；首字节锁存 \\
\end{longtable}

\section{入队与授权}

\code{afdx_tx_buf} 写完消息后，将 47 位描述符写入 \code{tx_desc}，再向引擎发出入队请求。引擎选中一条消息（或其中一片 IPv4 分片）后，将单帧/分片字段一次性交给组帧器；组帧器在帧开始发送时报告起始时刻，并在分片载荷读毕时回传下一片游标，见表~\ref{tab:ifgrant}。

\begin{longtable}{>{\raggedright\arraybackslash}p{4.4cm}>{\raggedleft\arraybackslash}p{1.8cm}>{\raggedright\arraybackslash}p{\dimexpr\linewidth-6.2cm-6\tabcolsep\relax}}
\caption{入队、分片授权和发送起始指示接口}\label{tab:ifgrant}\\
\longtablehead{3}{信号 & 位宽 & 含义}
\multicolumn{3}{l}{\textit{入队（\code{tx_buf} $\rightarrow$ \code{tx_engine}，请求应答）}} \\*
\code{enq_valid} & 1 & 保持有效，直至收到 \code{enq_ack} \\
\code{enq_vl}、\code{enq_head}、\code{enq_len} & 7、8、14 & 发送 VL 索引、消息首页号、整包字节数（$1\text{--}\MaxMessageBytes$） \\
\code{enq_ack}、\code{enq_result} & 1、3 & 单拍应答；0 接纳，1 VL 未启用，2 队列满，3 超长（4 无空闲页由 \code{tx_buf} 自行判定） \\
\code{tx_desc}[46:0] & 47 & \{\code{len}[13:0], \code{port_en}, \code{sport}[15:0], \code{dport}[15:0]\}，以首页号为地址 \\
\midrule
\multicolumn{3}{l}{\textit{授权（\code{tx_engine} $\rightarrow$ \code{tx_builder}，单深缓冲握手）}} \\*
\code{grant_valid}、\code{grant_ready} & 1、1 & \code{grant_ready} $=\neg$\code{pending}，即单帧待发槽位空闲 \\*
\code{grant_a}、\code{grant_b} & 1、1 & 两网启用位（配置字 0 的 [1]、[2]） \\*
\code{grant_timed}、\code{grant_time} & 1、32 & \code{grant_timed} 为 1 时，帧不得早于 \code{grant_time} 起始；为 0 时可立即发送 \\*
\code{grant_vlid} & 16 & 配置字 2 \\*
\code{grant_src}、\code{grant_dst} & 32、32 & 配置字 3、4 \\*
\code{grant_sport}、\code{grant_dport} & 16、16 & \code{port_en} 为 1 取消息自带端口，否则取配置字 5 \\*
\code{grant_len}、\code{grant_sn}、\code{grant_page} & 11、8、8 & 本帧载荷字节数、AFDX 帧序号、本片起始页号 \\*
\code{grant_head}、\code{grant_offset} & 8、9 & 消息首页号（首片保留至末片释放）、本片页内起始字节偏移 \\*
\code{grant_ip_offset}、\code{grant_message_len} & 13、14 & IPv4 8B 单位片偏移、整条应用消息总字节数 \\*
\code{grant_ident}、\code{grant_first}、\code{grant_more} & 16、1、1 & IPv4 包标识、首片标志（含 8B UDP 头）、\code{MF} 更多分片标志 \\*
\midrule
\multicolumn{3}{l}{\textit{发送起始与分片游标回传（\code{tx_builder} $\rightarrow$ \code{tx_engine}）}} \\*
\code{start_valid}、\code{t_start} & 1、32 & 单拍脉冲；\code{t_start} 为输出首个前导字节时的 \code{now} \\
\code{fragment_done}、\code{_vl}、\code{_page}、\code{_offset} & 1、7、8、9 & 非末片载荷读毕时回传下一待发分片的起始页号与页内偏移 \\
\end{longtable}

\section{组帧器与 MAC}

组帧器将目的 MAC 地址至 SN 字段的帧内容以字节流送至 MAC，由 MAC 控制字节接收时刻，见表~\ref{tab:ifs}。

\begin{longtable}{>{\raggedright\arraybackslash}p{4.4cm}>{\raggedleft\arraybackslash}p{1.8cm}>{\raggedright\arraybackslash}p{\dimexpr\linewidth-6.2cm-6\tabcolsep\relax}}
\caption{帧字节接口（\texttt{tx\_builder} $\leftrightarrow$ \texttt{mac\_tx}）}\label{tab:ifs}\\
\longtablehead{3}{信号 & 位宽 & 含义}
\code{s_valid}、\code{s_ready} & 1、1 & 字节握手；\code{s_ready} 在 \code{byte_en}$=1$ 且满足以下任一条件时为 1：前导码末拍；DATA 且 \code{last_byte}$=0$；DRAIN \\
\code{s_data_a}、\code{s_data_b} & 8、8 & 两网字节，仅第 11 字节不同 \\
\code{s_last} & 1 & 与 SN 字节同拍 \\
\code{s_en_a}、\code{s_en_b} & 1、1 & 起始拍锁存，选择本帧发往哪一网 \\
\code{tx_sof}、\code{tx_eof} & 1、1 & 首个前导字节、帧尾（FCS 发完）的单拍脉冲 \\
\code{tx_underrun} & 1 & 数据阶段上游未备好字节的下溢（Underrun）指示 \\
\end{longtable}

\section{帧描述与交付队列}

每网解析器为正确帧生成一份 221 位帧描述（兼容不分片单帧与 IPv4 分片帧）。引擎读取描述并完成 30 拍判定与重组槽位查询后，向解析器发出提交或回退指示；对于通过判定的帧，引擎向 $32\times 73$~bit 交付队列写入一项，见表~\ref{tab:ifdesc}。

\begin{longtable}{>{\raggedright\arraybackslash}p{4.4cm}>{\raggedleft\arraybackslash}p{1.8cm}>{\raggedright\arraybackslash}p{\dimexpr\linewidth-6.2cm-6\tabcolsep\relax}}
\caption{帧描述与交付队列项}\label{tab:ifdesc}\\
\longtablehead{3}{字段 & 位宽 & 含义}
\multicolumn{3}{l}{\textit{帧描述（\code{rx_parse} $\rightarrow$ \code{rx_engine}，单深缓冲握手，共 221 位，高位在前）}} \\*
\code{ip_src}、\code{ip_dst}、\code{ip_id} & 32、32、16 & IPv4 源/目的地址与 16 位分片标识（用于分片重组五元组与伪头） \\*
\code{frag_offset}、\code{mf} & 13、1 & IPv4 8B 单位片偏移与更多分片标志（均为 0 表示不分片单帧） \\*
\code{udp_len}、\code{udp_check} & 16、16 & 首片/单帧的 UDP 总长与校验和字段 \\*
\code{vlid}、\code{sn} & 16、8 & 目的 MAC 第 4--5 字节 VLID 与帧末 AFDX SN \\*
\code{sport}、\code{dport} & 16、16 & 首片/单帧提取的 UDP 源端口与目的端口 \\*
\code{n}、\code{start}、\code{arrival} & 11、12、16 & 本帧载荷字节数、环形缓冲起点与帧尾时刻的 \code{now_c} \\*
\code{desc_valid}、\code{desc_take} & 1、1 & 交接握手（按配置扇出分支复制） \\*
\code{resolve_valid}、\code{resolve_commit} & 1、1 & 引擎判定：提交则 \code{committed} 前进到本帧末，否则 \code{trial} 回退 \\*
\midrule
\multicolumn{3}{l}{\textit{交付队列项（\code{rx_engine} $\rightarrow$ \code{rx_deliver}，32 项 $\times$ 73~bit 块 RAM \code{rx_dq}，高位在前）}} \\*
\code{frag[72]}、\code{slot[71:68]} & 1、4 & 分片标志（0 单帧直通，1 分片重组）与预约的重组槽位号（\texttt{\{bank,row\}}） \\
\code{net[67]} & 1 & 0 为 A 网，1 为 B 网 \\
\code{start[66:55]}、\code{n[54:44]} & 12、11 & 环形缓冲起点和本帧载荷字节数 \\
\code{idx[43:32]} & 12 & 接收 VL 索引 \\
\code{sport[31:16]}、\code{dport[15:0]} & 16、16 & UDP 源端口与目的端口（共 73 位） \\
\end{longtable}

核心交付口（重组模块输出 \code{d_*}）的字段与 \code{afdx_es} 应用接收口对应（目的端口在内部记为 \code{d_dport}，在核心记为 \code{app_rx_port}）。服务通道将这些信号同时广播给应用和三个服务。各服务在首字节到达当拍，根据 UDP 目的端口产生组合匹配信号 \code{svc_rx_claim}；匹配结果寄存并保持至末字节。端口匹配的消息由对应服务接收；其余应用消息在顶层 \code{AFDX_top} 中经 \code{afdx_app_rx} 加上 2 字节大端 VLID 前缀后自 \code{rx_data}/\code{rx_valid} 输出。

\section{服务发送总线}

三个服务的发送口打包成一组总线接到 \code{afdx_svc}，按 RTC、SNMP、615A 从低位到高位拼接，见表~\ref{tab:ifsvc}。各服务按相同的字节流握手约定驱动对应位段。

\begin{longtable}{>{\raggedright\arraybackslash}p{4.4cm}>{\raggedleft\arraybackslash}p{1.8cm}>{\raggedright\arraybackslash}p{\dimexpr\linewidth-6.2cm-6\tabcolsep\relax}}
\caption{服务发送总线 \texttt{svc\_tx\_*}}\label{tab:ifsvc}\\
\longtablehead{3}{信号 & 位宽 & 拼接}
\code{svc_tx_valid}、\code{_ready}、\code{_last}、\code{_port_en} & 各 $3\times1$ & [0] RTC，[1] SNMP，[2] 615A \\
\code{svc_tx_data} & $3\times8$ & [7:0] RTC，[15:8] SNMP，[23:16] 615A \\
\code{svc_tx_vl} & $3\times7$ & [6:0]、[13:7]、[20:14] \\
\code{svc_tx_sport}、\code{svc_tx_dport} & 各 $3\times16$ & [15:0]、[31:16]、[47:32] \\
\end{longtable}

\chapter{术语与参考文献}

\section{术语}\label{app:terms}

\begin{table}\centering\small
\caption{本文专用术语}\label{tab:terms}
\begin{tabularx}{\linewidth}{>{\raggedright\arraybackslash}p{3.4cm}L}
\toprule
术语 & 含义 \\
\midrule
拍 & 一个 125~MHz 时钟周期，8~ns \\
PSN & 同一网上同一 VL 上一帧的序列号 \\
预写入 / 提交 / 回退 & 载荷先写入环形缓冲；通过检查后更新提交写指针，拒绝时回退预写入指针 \\
候选 & 调度引擎选出、等待授权给组帧器的 VL \\
\code{pend_vl} & 已授权但尚未开始发送的 VL \\
\code{cfg_ok} & 引擎在本拍不发出配置表读的指示 \\
有效副本 & 帧格式、FCS 与协议头检查通过，且满足接收配置和已启用的完整性检查的单网帧副本 \\
页面释放 & 将页号送回空闲页池（经页面回收口 \code{ret1_*} 或 \code{ret2_*}）\\
端口匹配 & 服务根据 UDP 目的端口确定是否接收当前消息（信号 \code{svc_rx_claim}）\\
表初始化 & 复位后逐项清零收发配置表与状态表（\code{rx_map}、\code{rx_cfg}、\code{rx_state}、\code{tx_cfg}、\code{tx_state}）的过程 \\
线路占用 & 一帧在线路上的占用时长 $C(n)$，单位为拍 \\
Server / Client & SNTP 角色：Server 端应答请求并回显时间戳；Client 端周期发送请求并执行时钟同步校准 \\
应答 & 请求应答式握手的确认脉冲（\code{ack}），用于确认接口请求，与协议响应报文不同 \\
计时单位 & \code{now_c} 的单位，32.768~µs \\
发送抖动 & 帧起始相对 BAG 允许时刻的偏差 \\
起始指示 / 起始确认 & 组帧器在 \code{tx_sof} 有效拍经 \code{start_valid}/\code{t_start} 发送起始指示（TX Start Indication）；发送引擎据此完成发送起始确认（TX Start Confirm）并写回更新后的 \code{next_time} \\
\bottomrule
\end{tabularx}
\end{table}

\section{参考文献}\label{app:ref}

\begingroup
\small\linespread{1.08}\selectfont
\begin{enumerate}[label={[\arabic*]},leftmargin=2.5em]
  \item ARINC Specification 664 Part 7: Avionics Full Duplex Switched Ethernet Network, 2005-06-27.
  \item ARINC Report 615A-2: Software Data Loader Using Ethernet Interface, 2002-05-10.
  \item ARINC Report 665-1: Loadable Software Standards, 2001-01-12.
  \item RFC 1350: The TFTP Protocol (Revision 2), 1992.
  \item RFC 2347: TFTP Option Extension, 1998.
  \item RFC 1157: A Simple Network Management Protocol (SNMP), 1990.
  \item RFC 3416: Version 2 of the Protocol Operations for the Simple Network Management Protocol (SNMP), 2002.
  \item RFC 3418: Management Information Base (MIB) for the Simple Network Management Protocol (SNMP), 2002.
  \item RFC 4330: Simple Network Time Protocol (SNTP) Version 4 for IPv4, IPv6 and OSI, 2006.
  \item IEEE Std 802.3: IEEE Standard for Ethernet, 2018.
  \item 中科亿海微. 全国大学生嵌入式芯片与系统设计竞赛 FPGA 创新设计赛道选题指南.
\end{enumerate}
\endgroup
